\begin{table}[!htbp]
\centering
\caption{Modelo principal, 3ª série EM em Matemática: ATT por período relativo}
\label{tab:ap_event_base_em3_mat}
\vspace{-0.15cm}
\scriptsize
\setlength{\tabcolsep}{4pt}
\renewcommand{\arraystretch}{0.92}
\resizebox{\textwidth}{!}{%
\begin{tabular}{lcccc}
\toprule
Período relativo & TWFE & CS21 & DR-DiD & S\&X aprox. \\
\midrule
-4 & 0.064 (0.029) & 0.059 (0.040) & 0.056 (0.041) & 0.075 (0.077) \\
-3 & 0.059 (0.034) & 0.038 (0.035) & 0.031 (0.035) & 0.091 (0.049) \\
-2 & 0.051 (0.029) & 0.047 (0.029) & 0.061 (0.031) & 0.087 (0.041) \\
-1 & 0.000 & 0.000 & 0.000 & 0.000 \\
0 & 0.346 (0.037) & 0.332 (0.037) & 0.349 (0.039) & 0.410 (0.065) \\
1 & 0.485 (0.039) & 0.469 (0.042) & 0.488 (0.043) & 0.552 (0.067) \\
2 & 0.665 (0.041) & 0.651 (0.045) & 0.667 (0.045) & 0.750 (0.081) \\
3 & 0.717 (0.046) & 0.736 (0.053) & 0.767 (0.051) & 0.776 (0.094) \\
4 & 0.806 (0.047) & 0.840 (0.061) & 0.865 (0.059) & 0.887 (0.103) \\
\bottomrule
\end{tabular}
}
\begin{flushleft}
\footnotesize Nota: cada célula reporta ATT e erro-padrão entre parênteses, em desvios-padrão da proficiência anual. O período \(k=-1\) é normalizado como linha de base; por isso aparece sem erro-padrão.
\end{flushleft}
\end{table}
